% Generated draft: requires mathematical and source review.
\documentclass[11pt]{article}
\usepackage[margin=25mm]{geometry}
\usepackage{fontspec}
\setmainfont{Latin Modern Roman}
\usepackage{amsmath,amssymb,mathtools,unicode-math}
\setmathfont{Latin Modern Math}
\usepackage{enumitem,needspace,xurl,hyperref}
\setlength{\parindent}{0pt}
\setlength{\parskip}{5pt}
\setlength{\emergencystretch}{3em}
\newcommand{\N}{\mathbb{N}}
\newcommand{\Z}{\mathbb{Z}}
\newcommand{\Q}{\mathbb{Q}}
\newcommand{\R}{\mathbb{R}}
\newcommand{\C}{\mathbb{C}}
\newcommand{\F}{\mathbb{F}}
\newcommand{\A}{\mathbb{A}}
\newcommand{\PP}{\mathbb{P}}
\newcommand{\E}{\mathbb{E}}
\newcommand{\eps}{\varepsilon}
\newcommand{\dd}{\,\mathrm{d}}
\begin{document}
\section*{UnsolvedMath ID: 701618}
\textbf{Source problem code:} catalogue entry
\subsection{Definitions and mathematical statement}
A topological group is \emph{tdlc} if it is totally disconnected and locally compact, and is \emph{tdlcsc} if it is additionally second countable. A group is \emph{compactly generated} if some compact subset generates it as an abstract group. A topological group is \emph{topologically simple} if its only closed normal subgroups are the trivial subgroup and the whole group.

Interpret \emph{elementary} in the standard tdlcsc sense: let \(\mathcal E\) be the smallest class of tdlcsc groups, closed under topological isomorphism, containing all countable discrete groups and all second-countable profinite groups, and closed under group extensions and countable increasing unions of open subgroups. Here closure under extensions means that \(L\in\mathcal E\) whenever \(L\) is tdlcsc and has a closed normal subgroup \(M\) with \(M,L/M\in\mathcal E\); closure under unions means that \(L\in\mathcal E\) whenever \(L=\bigcup_{n\geq1}L_n\), where \(L_n\leq L_{n+1}\) are open subgroups belonging to \(\mathcal E\). A tdlcsc group is \emph{non-elementary} if it does not belong to \(\mathcal E\). The supplied record does not explicitly define this terminology, so this interpretation requires confirmation against the source.

Let
\[
\mathscr S=\{S: S\text{ is a nondiscrete, compactly generated, topologically simple tdlc group}\}.
\]
The question is whether, for every non-elementary tdlcsc group \(G\), there exist a closed subgroup \(H\leq G\), compactly generated with its subspace topology, and a closed normal subgroup \(N\trianglelefteq H\) such that
\[
H/N\in\mathscr S,
\]
where \(H/N\) carries the quotient topology. This interprets the source's phrase ``continuous quotient'' as a Hausdorff topological group quotient.
\subsection{Short English statement}
Does every non-elementary totally disconnected, locally compact, second-countable group contain a compactly generated closed subgroup with a nondiscrete, compactly generated, topologically simple quotient? The meanings of elementary and continuous quotient are made explicit above, subject to source confirmation.
\subsection{Sources}
\begin{itemize}
\item Local UnsolvedMath registry, supplied complete record, permanent ID \(701618\), exact source problem code \texttt{catalogue entry}, title \emph{Kourovka Notebook Problem 19.70}. The registry labels the problem open; this is supplied metadata, not an independently verified status assessment.
\item Supplied source attribution: \emph{The Kourovka Notebook (Unsolved Problems in Group Theory)}, 19th issue (2018), original problem 19.70; the background also identifies the 21st-issue edition (2026), \texttt{arXiv:1401.0300v48}. The registry describes the collection as edited by E. I. Khukhro and V. D. Mazurov. These bibliographic details are supplied and have not been independently checked.
\item The statement credits P. Wesolek and P.-E. Caprace. The registry's proposer field and background instead credit C. Reid, with proposed year 2018. The supplied record does not explain these respective attribution roles.
\end{itemize}
\end{document}
